% =============================================================================
% Day 1 Course Overview — ML Systems Lecture Slides
% =============================================================================
% This deck is presented on the FIRST DAY of class to introduce students
% to the course identity, philosophy, analytical frameworks, and logistics.
%
% IMAGES NEEDED (generate SVGs, convert to PDF):
%   - course-roadmap.pdf       - deployment-spectrum.pdf
%   - frameworks-overview.pdf  - lighthouse-models.pdf
%   - five-percent-overview.pdf
% =============================================================================
\documentclass[aspectratio=169, 12pt]{beamer}
\usepackage{../../assets/beamerthememlsys}

\mlsyssetup{
  volume       = {Volume I},
  chapter      = {Course Overview},
  logo         = {../../assets/img/logo-mlsysbook.png},
  instlogo     = {../../assets/img/logo-harvard.png},
  chaptertitle = {Introduction to ML Systems},
}

% --- Fonts ---
\usepackage{fontspec}
\setsansfont{Helvetica Neue}[
  BoldFont={Helvetica Neue Bold},
  ItalicFont={Helvetica Neue Italic},
  BoldItalicFont={Helvetica Neue Bold Italic},
]
\setmonofont{JetBrains Mono}[Scale=0.85]

% --- Packages ---
\usepackage{booktabs}
\usepackage{amsmath}

% --- Image paths ---
\graphicspath{
  {images/}
}

% --- Chapter-specific macros ---
\newcommand{\DAM}{D\raisebox{0.08em}{\tiny$\bullet$}A\raisebox{0.08em}{\tiny$\bullet$}M}

% --- Helper: safe image include ---
\newcommand{\safeimg}[2][width=\textwidth,keepaspectratio]{%
  \IfFileExists{images/#2}{\includegraphics[#1]{#2}}{%
    \IfFileExists{#2}{\includegraphics[#1]{#2}}{%
      \fbox{\parbox[c][2.5cm][c]{0.85\linewidth}{\centering\footnotesize\textcolor{midgray}{[Missing image]}}}%
    }%
  }%
}


\begin{frame}{Visual Language}
\note{
% -- NARRATE: Walk through each card. ``Blue means compute---anytime you see
% blue in a diagram, think GPU ops, matrix multiplies, forward pass. Green is
% data flow---memory, caches, healthy paths. Orange is routing---schedulers,
% load balancers. Red is cost, error, or bottleneck.'' Point to each card as
% you name it.
%
% -- FLEX: [CORE] Show on Day 1 and again briefly on Day 2.
% IF SHORT: Display the slide but do not narrate---students can read the cards.
}

\small
Throughout this course, colors carry meaning:

\vspace{0.3cm}
\begin{columns}[T]
  \begin{column}{0.45\textwidth}
    \begin{mlsyscard}{computestroke}
    \textbf{Blue} --- Compute / Processing\\
    {\footnotesize GPU ops, forward/backward pass, inference}
    \end{mlsyscard}
    \vspace{0.15cm}
    \begin{mlsyscard}{datastroke}
    \textbf{Green} --- Data / Memory\\
    {\footnotesize Data flow, caches, healthy paths}
    \end{mlsyscard}
  \end{column}
  \begin{column}{0.45\textwidth}
    \begin{mlsyscard}{routingstroke}
    \textbf{Orange} --- Routing / Scheduling\\
    {\footnotesize Load balancers, batch windows}
    \end{mlsyscard}
    \vspace{0.15cm}
    \begin{mlsyscard}{errorstroke}
    \textbf{Red} --- Error / Cost / Bottleneck\\
    {\footnotesize Loss, decode phase, waste}
    \end{mlsyscard}
  \end{column}
\end{columns}
\end{frame}

% --- Section count for navigation (must match actual \section{} count) ---
\setcounter{mlsystotalsections}{7}

\title{Course Overview}
\author{Vijay Janapa Reddi}
\institute{Harvard University}
\date{}

\begin{document}

% =============================================================================
% TITLE SLIDE
% =============================================================================
\mlsystitle{Course Overview}{The Physics of AI Engineering}{}

% =============================================================================
% WELCOME & COURSE IDENTITY
% =============================================================================

\begin{frame}{Welcome}
\note{
% -- NARRATE: ``Welcome. Raise your hand if you have trained a model.'' (most
% hands go up) ``Now keep your hand up if you have deployed one to production.''
% (most hands drop) ``That gap---training vs.\ shipping---is this entire course.
% This is not an ML algorithms class. This is about the systems that make ML
% work: memory, bandwidth, power, latency, and the physics behind every design
% decision.''
%
% -- FLEX: [CORE] Sets the emotional contract for the semester.
% IF SHORT: Cut the hand-raise; just state the training-deployment gap directly.
}

\centering
\vspace{0.8cm}
{\Large\bfseries Introduction to Machine Learning Systems}

\vspace{0.4cm}
{\large Volume I: Foundations for Single-Machine ML Systems}

\vspace{0.6cm}
{\normalsize\textcolor{crimson}{\textbf{``The Physics of AI Engineering''}}}

\vspace{0.6cm}
{\small AI is not magic --- it is infrastructure, and infrastructure has laws.}

\end{frame}

% =============================================================================
\section{Why This Course}
% =============================================================================

\begin{frame}{The Gap Between ML Research and Production}
\note{
% -- LINK: Students just heard ``this is a systems course, not an algorithms
% course.'' This slide gives the quantitative evidence for WHY.
%
% -- NARRATE: ``60 to 85 percent of ML projects never reach production. Let
% that sink in. Look at this table---research uses a static dataset, a single
% GPU, and optimizes one metric. Production uses a shifting data stream, a
% fleet of heterogeneous hardware, and must hit accuracy AND latency AND cost
% targets simultaneously. That gap is not fixed by a better optimizer.''
%
% -- ENGAGE: ``Why do you think most ML projects fail? Write one reason.''
% Give 30 seconds. Cold-call 2 students.
% Expected: ``data quality,'' ``hardware limits.'' Surprise answer: systems.
%
% -- WARN: Students assume failures are algorithmic (``bad model''). Correct
% framing: the bottleneck is infrastructure---data pipelines, serving, monitoring.
% IF STUCK: Point to the ``90\% of time goes to data + infrastructure'' bullet.
%
% -- FLEX: [CORE] This slide motivates the entire semester.
% IF AHEAD: ``What percentage of engineering time goes to the model itself?''
}

\small
\begin{columns}[T]
  \begin{column}{0.55\textwidth}
    A Kaggle notebook is not a production system:

    \vspace{0.2cm}
    \begin{itemize}\setlength\itemsep{1pt}
      \item \textbf{60--85\%} of ML projects fail to reach production
      \item 90\% of engineering time goes to data + infrastructure
      \item Model code is $\approx$\textbf{5\%} of a production system
    \end{itemize}

    \vspace{0.2cm}
    \begin{mlsyscard}{errorstroke}
    The bottleneck is \alert{systems}, not algorithms.
    \end{mlsyscard}
  \end{column}
  \begin{column}{0.42\textwidth}
    \vspace{0.3cm}
    \renewcommand{\arraystretch}{1.25}
    {\footnotesize
    \begin{tabular}{@{}lll@{}}
      \toprule
      & \textbf{Research} & \textbf{Production} \\
      \midrule
      Data       & Static set   & Shifting stream \\
      Hardware   & Given GPU    & Fleet of tiers \\
      Success    & Accuracy     & Accuracy + SLA \\
      Lifespan   & One paper    & Years \\
      \bottomrule
    \end{tabular}
    }
  \end{column}
\end{columns}
\mlsyscite{Paleyes et al., ACM Computing Surveys 2022}

\end{frame}

\begin{frame}{The 5\% Problem}
\note{
% -- LINK: The previous slide said ``the bottleneck is systems.'' This diagram
% shows exactly what those systems look like.
%
% -- NARRATE: Point to the tiny center box: ``That is the ML model code---about
% 5 percent. Everything around it---data pipelines, feature stores, serving
% infrastructure, monitoring, configuration---is what this course teaches.
% Sculley et al.\ called this `hidden technical debt.' ''
%
% -- ENGAGE: ``Look at the diagram. What is the biggest box?'' Give 15 seconds.
% Expected answer: data collection or configuration. Both are valid---the point
% is that neither is the model.
%
% -- WARN: Students equate ``ML'' with ``the model.'' Correct framing: the
% model is 5\%; the other 95\% is systems engineering that determines whether
% the model ever reaches a user.
%
% -- FLEX: [CORE] Foundational mental model for the course.
% IF SHORT: Skip the question; just narrate the diagram for 90 seconds.
}

% --- Layout: FULL-WIDTH IMAGE ---
\centering
\safeimg[width=0.88\textwidth,height=4.8cm,keepaspectratio]{five-percent-overview.pdf}

\vspace{0.1cm}
{\small This course teaches the \textbf{95\%} --- the systems engineering that makes ML work.}
\mlsyscite{Sculley et al., NeurIPS 2015}

\end{frame}

% =============================================================================
\section{The Core Thesis}
% =============================================================================

\mlsysfocus{The Core Thesis}{%
Constraints drive architecture.\\[0.3cm]
You don't choose a Transformer because it's trendy ---\\
you choose it because of how it parallelizes on real silicon.%
}

\begin{frame}{AI Is Infrastructure}
\note{
% -- LINK: The Core Thesis focus slide just said ``constraints drive
% architecture.'' This slide names the four specific physical constraints.
%
% -- NARRATE: Walk through bullets top to bottom. ``Memory bandwidth limits
% how fast data reaches the processor---this is why GPT-4 inference is slow
% even on powerful GPUs. Power budget limits where a model can run---a phone
% cannot sustain 700 watts. Speed of light limits latency---a self-driving car
% cannot wait 50 ms for a cloud round trip. Thermodynamics limits compute per
% rack---you cannot cool infinite GPUs in a data center.''
% ANALOGY: ``Physics is to ML systems what gravity is to bridges. You can
% build creative bridges, but none of them ignore gravity.''
%
% -- ENGAGE: ``Why can't we run GPT-4 on a smartphone?'' Cold-call one student.
% Expected: ``not enough memory.'' Deepen: ``How much memory does it need?
% 3.6 TB at FP32. A phone has 8 GB. That is a 450x gap---physics, not software.''
%
% -- WARN: Students think hardware limitations are temporary (``next year's chip
% will fix it''). Correct framing: these are physical laws, not engineering gaps.
% Memory bandwidth grows ~20\%/yr; compute demand grows 7x faster than Moore's Law.
%
% -- FLEX: [CORE] Philosophical anchor for the course.
% IF AHEAD: ``Which constraint is hardest to overcome with engineering?''
}

\small
\begin{columns}[T]
  \begin{column}{0.55\textwidth}
    \textbf{Physical laws govern ML systems:}

    \vspace{0.1cm}
    \begin{itemize}\setlength\itemsep{0pt}
      \item \textcolor{computestroke}{\textbf{Memory bandwidth}} limits how fast data reaches the processor
      \item \textcolor{routingstroke}{\textbf{Power budget}} limits where a model can run
      \item \textcolor{errorstroke}{\textbf{Speed of light}} limits latency over distance
      \item \textcolor{datastroke}{\textbf{Thermodynamics}} limits compute per rack
    \end{itemize}

    \vspace{0.05cm}
    \mlsysconcept{Implication:}{Physics, not elegance, determines what ships.}
  \end{column}
  \begin{column}{0.42\textwidth}
    \vspace{0.1cm}
    \begin{mlsyscard}{crimson}
    {\footnotesize \textbf{Example:} GPT-4 costs $\sim$\$100M to train --- moving 1.8T parameters through memory thousands of times costs that much in electricity and silicon time.}
    \end{mlsyscard}

    \vspace{0.1cm}
    \begin{mlsyscard}{datastroke}
    {\footnotesize \textbf{Example:} A self-driving car cannot call the cloud --- speed of light adds 50+ ms. Physics forces the model onto the edge.}
    \end{mlsyscard}
  \end{column}
\end{columns}

\end{frame}

% =============================================================================
\section{Three Frameworks}
% =============================================================================

\begin{frame}{Three Analytical Frameworks}
\note{
% -- LINK: Students now know that physical constraints drive architecture.
% These three frameworks are the tools for reasoning about those constraints.
%
% -- NARRATE: Point to each framework in the diagram. ``D-A-M tells you WHERE
% the bottleneck is. The Iron Law tells you HOW LONG an operation takes. The
% Degradation Equation tells you WHEN a model will fail. These three tools
% recur in every single chapter.''
%
% -- ENGAGE: ``What question does each framework answer? Write one word per
% framework.'' 30 seconds. Cold-call one student.
% Expected: Where/How long/When---or close variants.
%
% -- WARN: Students will try to memorize the equations without understanding
% what question each answers. Correct framing: frameworks are diagnostic
% tools, not formulas to plug numbers into.
%
% -- FLEX: [CORE] Preview only---do not go deep. Chapter 1 covers each.
% IF SHORT: Just name the three frameworks and move on (60 seconds).
}

% --- Layout: FULL-WIDTH IMAGE ---
\centering
\safeimg[width=0.92\textwidth,height=4.8cm,keepaspectratio]{frameworks-overview.pdf}

\vspace{0.1cm}
{\small Three tools you will use in \textbf{every chapter} of this course.}

\end{frame}

\begin{frame}{The \DAM{} Taxonomy}
\note{
% -- LINK: The previous diagram named D-A-M as one of three frameworks.
% Now we unpack it briefly.
%
% -- NARRATE: ``Every ML system sits at the intersection of three axes.
% Data: how much, how fast can we move it. Algorithm: how many operations,
% what parallelism pattern. Machine: what silicon, what memory hierarchy.
% The diagnostic question is always the same: which axis is the bottleneck?
% Optimizing one axis shifts pressure to the others---they are coupled.''
% Point to the table: ``Notice the units. These become the Iron Law variables.''
%
% -- ENGAGE: ``If you double the training dataset, which other axis feels
% the pressure?'' Expected: Machine (need more bandwidth or compute time).
%
% -- WARN: Students treat the three axes as independent knobs. Correct framing:
% they are interdependent---doubling data volume requires proportionally more
% bandwidth or longer training time.
%
% -- FLEX: [CORE] But keep it brief---Chapter 1 goes deep.
% IF SHORT: Show the slide for 60 seconds, skip the engage question.
}

\small
\begin{columns}[T]
  \begin{column}{0.55\textwidth}
    Every ML system has three interdependent axes:

    \vspace{0.2cm}
    \textcolor{datastroke}{\textbf{Data}} --- Volume, quality, distribution\\
    {\footnotesize How much data, how fast can we move it?}

    \vspace{0.1cm}
    \textcolor{computestroke}{\textbf{Algorithm}} --- Architecture, operation count\\
    {\footnotesize How many FLOPs, what parallelism pattern?}

    \vspace{0.1cm}
    \textcolor{routingstroke}{\textbf{Machine}} --- Hardware throughput, memory\\
    {\footnotesize What silicon, what memory hierarchy?}

    \vspace{0.15cm}
    \mlsysalert{Key:}{Optimizing one axis shifts pressure to another.}
  \end{column}
  \begin{column}{0.42\textwidth}
    \vspace{0.3cm}
    \renewcommand{\arraystretch}{1.2}
    {\footnotesize
    \begin{tabular}{@{}lll@{}}
      \toprule
      \textbf{Axis} & \textbf{Variable} & \textbf{Unit} \\
      \midrule
      Data      & $D_{\text{vol}}$, $BW$ & GB, GB/s \\
      Algorithm & $O$ & FLOPs \\
      Machine   & $R_{\text{peak}}$ & FLOP/s \\
      \bottomrule
    \end{tabular}
    }

    \vspace{0.2cm}
    \begin{mlsyscard}{crimson}
    {\footnotesize \textbf{Diagnostic:} ``Which axis is the bottleneck?'' answers most ML systems design questions.}
    \end{mlsyscard}
  \end{column}
\end{columns}

\end{frame}

\begin{frame}{The Iron Law of ML Systems}
\note{
% -- LINK: D-A-M tells you where the bottleneck is. The Iron Law quantifies
% how long each axis takes in seconds.
%
% -- NARRATE: Point to each term. ``Data term: bytes divided by bandwidth
% gives seconds. Compute term: FLOPs divided by peak rate times efficiency
% gives seconds. Overhead: orchestration tax, also seconds. You add three
% times and the slowest one dominates. This is dimensional analysis---if
% your units do not resolve to seconds, the equation is wrong.''
%
% -- ENGAGE: ``For a phone camera app classifying a photo, which term
% dominates?'' Give 20 seconds. Expected: Data term (reading the image from
% memory) or Overhead (framework launch cost). Accept either with reasoning.
%
% -- WARN: Students try to add FLOPs to bytes. Correct framing: every term
% must resolve to seconds before you can compare or add them.
%
% -- FLEX: [CORE] Preview only---worked examples come in Chapter 1.
% IF SHORT: State the equation, emphasize ``slowest term dominates,'' move on.
}

\small
$$T_{\text{total}} = \underbrace{\dfrac{D_{\text{vol}}}{BW}}_{\text{Data}} +
\underbrace{\dfrac{O}{R_{\text{peak}} \cdot \eta}}_{\text{Compute}} +
\underbrace{L_{\text{lat}}}_{\text{Overhead}}$$

\vspace{0.1cm}
\begin{columns}[T]
  \begin{column}{0.30\textwidth}
    \centering
    \textcolor{datastroke}{\textbf{Data}} $\frac{D_{\text{vol}}}{BW}$\\[0.05cm]
    {\scriptsize Bytes / (Bytes/s) = \textbf{s}}
  \end{column}
  \begin{column}{0.30\textwidth}
    \centering
    \textcolor{computestroke}{\textbf{Compute}} $\frac{O}{R_p \cdot \eta}$\\[0.05cm]
    {\scriptsize FLOPs / (FLOPs/s) = \textbf{s}}
  \end{column}
  \begin{column}{0.30\textwidth}
    \centering
    \textcolor{errorstroke}{\textbf{Overhead}} $L_{\text{lat}}$\\[0.05cm]
    {\scriptsize Orchestration tax = \textbf{s}}
  \end{column}
\end{columns}

\vspace{0.15cm}
\centering
\mlsysconcept{Key insight:}{Every term resolves to \textbf{seconds}. The \textbf{slowest term dominates}.}

\end{frame}

\begin{frame}{The Degradation Equation}
\note{
% -- LINK: The Iron Law measures performance at a point in time. The
% Degradation Equation measures how performance decays over time.
%
% -- NARRATE: ``Accuracy at time t equals initial accuracy minus alpha times
% the distribution distance. Alpha is how sensitive the model is to drift.
% Delta measures how far the live data has drifted from training data.
% Look at the example: a recommendation system starts at 85\% and drops to
% 79.2\% in 6 months. No code changed. No bugs. The world changed.
% The engineering response: set a retraining trigger at a threshold.''
%
% -- ENGAGE: ``How would you know your model is getting worse if no code
% changed and no one filed a bug?'' Give 20 seconds. Cold-call one student.
% Expected: monitoring accuracy metrics over time. Deepen: ``What if you
% do not have ground truth labels in real time?''
%
% -- WARN: Students assume ``no bugs = working correctly.'' Correct framing:
% ML systems degrade through data drift even when code is untouched.
%
% -- FLEX: [CORE] Third framework preview.
% IF SHORT: Show equation, state the rec-system example, move on.
}

\small
$$A(t) = A_0 - \alpha \cdot \Delta(P_{\text{train}},\; P_{\text{live}}(t))$$

\vspace{0.15cm}
\begin{columns}[T]
  \begin{column}{0.50\textwidth}
    {\footnotesize
    \begin{tabular}{@{}ll@{}}
      $A(t)$    & Accuracy at time $t$ \\
      $A_0$     & Accuracy at deployment \\
      $\alpha$  & Sensitivity to drift \\
      $\Delta$  & Distribution distance \\
    \end{tabular}
    }

    \vspace{0.15cm}
    \mlsysalert{Key:}{No code changed. The \emph{world} changed.}
  \end{column}
  \begin{column}{0.46\textwidth}
    \begin{mlsyscard}{errorstroke}
    \textbf{Example:}\\[0.1cm]
    {\footnotesize Recommendation system: 85\% accuracy at launch $\to$ 79.2\% after 6 months. No bugs --- just user behavior shift.\\[0.1cm]
    \textbf{Engineering response:} Set retraining triggers when $A(t)$ crosses a threshold.}
    \end{mlsyscard}
  \end{column}
\end{columns}

\end{frame}

% --- ACTIVE LEARNING 1: Predict ---
\begin{frame}{Predict: Which Constraint?}
\note{[2 min] Prediction exercise. Give students 60 seconds. The point is to
prime their intuition before we introduce the lighthouse models.
Do NOT reveal the answer --- just let them think.}

\centering
\vspace{0.8cm}
{\Large\bfseries Think--Write--Share}

\vspace{0.6cm}
{\large A voice assistant must respond in under 200 ms.\\[0.2cm]
\alert{What physical constraint dominates?}}

\vspace{0.4cm}
\begin{columns}[c]
  \begin{column}{0.22\textwidth}\centering\small Memory\\ Bandwidth\end{column}
  \begin{column}{0.22\textwidth}\centering\small Compute\\ Throughput\end{column}
  \begin{column}{0.22\textwidth}\centering\small Network\\ Latency\end{column}
  \begin{column}{0.22\textwidth}\centering\small Power\\ Budget\end{column}
\end{columns}

\vspace{0.5cm}
{\small\textcolor{lightgray}{Write your answer. Then turn to a neighbor.}}

\end{frame}

% =============================================================================
\section{Lighthouse Models}
% =============================================================================

\begin{frame}{Five Lighthouse Models}
\note{[3 min] These five models are the recurring cast of characters.
Each isolates a distinct systems bottleneck. They recur in every chapter.
Walk through the diagram. Ask: ``Why do we need five, not one?''
Answer: because different models stress different parts of the system.}

% --- Layout: FULL-WIDTH IMAGE ---
\centering
\safeimg[width=0.92\textwidth,height=4.8cm,keepaspectratio]{lighthouse-models.pdf}

\vspace{0.1cm}
{\small Five models, five bottlenecks --- your \textbf{systems detectives} for the semester.}

\end{frame}

\begin{frame}{Lighthouse Models: Quick Reference}
\note{[2 min] Compact table for reference. Students will see these numbers
repeatedly. Emphasize: same Iron Law equation, different dominant term.}

\footnotesize
\renewcommand{\arraystretch}{1.2}
\begin{tabular}{@{}lp{2cm}lp{4cm}@{}}
  \toprule
  \textbf{Lighthouse} & \textbf{Bottleneck} & \textbf{Iron Law} & \textbf{Systems Question} \\
  \midrule
  \textbf{ResNet-50}    & \textcolor{computestroke}{Compute}  & $O/R_p$ & Is arithmetic fast enough? \\
  \textbf{GPT-2}        & \textcolor{datastroke}{Memory-BW}   & $D/BW$  & Can memory feed the accelerator? \\
  \textbf{MobileNetV2}  & \textcolor{routingstroke}{Latency}  & $L_{\text{lat}}$ & Can it meet real-time deadlines? \\
  \textbf{DLRM}         & \textcolor{datastroke}{Data I/O}    & $D/BW$  & Can we move embeddings fast enough? \\
  \textbf{KWS}          & \textcolor{errorstroke}{Power}      & $P \times T$ & Can it run on a coin-cell battery? \\
  \bottomrule
\end{tabular}

\vspace{0.15cm}
{\footnotesize Same equation, different dominant term. Each lighthouse tests a different constraint regime.}

\end{frame}

% =============================================================================
\section{Course Roadmap}
% =============================================================================

\begin{frame}{Sixteen Chapters, Four Parts}
\note{[3 min] Walk through the roadmap. Emphasize the progression:
Understand $\to$ Build $\to$ Optimize $\to$ Ship. Each part builds on the
previous --- skip nothing. The recurring themes (D-A-M, Iron Law, lighthouses)
tie everything together.}

% --- Layout: FULL-WIDTH IMAGE ---
\centering
\safeimg[width=0.92\textwidth,height=4.8cm,keepaspectratio]{course-roadmap.pdf}

\vspace{0.1cm}
{\small \textbf{Understand} $\to$ \textbf{Build} $\to$ \textbf{Optimize} $\to$ \textbf{Ship} --- each part builds on the previous.}

\end{frame}

\begin{frame}{The Deployment Spectrum}
\note{[2 min] Nine orders of magnitude separate Cloud and TinyML.
You cannot ``shrink'' a cloud model to the edge --- each tier requires
fundamentally different engineering. This is the physical reality
that drives every design decision in the course.}

% --- Layout: FULL-WIDTH IMAGE ---
\centering
\safeimg[width=0.92\textwidth,height=4.8cm,keepaspectratio]{deployment-spectrum.pdf}

\vspace{0.1cm}
{\small $10^5\!\times$ memory, $10^6\!\times$ compute --- you cannot ``just shrink'' a cloud model.}

\end{frame}

\begin{frame}{Quick Check}
\note{[0.5 min] Micro-retrieval. Give students 30 seconds to recall.
This distributed practice improves retention (Roediger \& Karpicke, 2006).}

\centering
\Large\bfseries Quick check --- no peeking.\\[0.5cm]
\normalsize Which physical constraint prevents running GPT-4 on a smartphone?\\[0.3cm]
{\small 30 seconds --- then we continue.}
\end{frame}



% --- ACTIVE LEARNING 2: Discussion ---
\begin{frame}{Discussion: Where Does It Run?}
\note{[3 min] Turn-and-talk. Students discuss in pairs for 90 seconds.
This connects deployment spectrum to real decisions. Cold-call 2--3 pairs.
Multiple right answers exist --- the point is the reasoning, not the answer.}

\centering
\vspace{0.8cm}
{\large\bfseries Turn and Talk \textcolor{midgray}{(90 seconds)}}

\vspace{0.6cm}
{\large A hospital wants to deploy a model that detects\\
pneumonia from chest X-rays in real time.\\[0.3cm]
\alert{Where on the deployment spectrum should it run?}}

\vspace{0.5cm}
\begin{columns}[c]
  \begin{column}{0.22\textwidth}\centering\small\textcolor{computestroke}{\textbf{Cloud}}\\700 W\end{column}
  \begin{column}{0.22\textwidth}\centering\small\textcolor{datastroke}{\textbf{Edge}}\\15--40 W\end{column}
  \begin{column}{0.22\textwidth}\centering\small\textcolor{routingstroke}{\textbf{Mobile}}\\3--5 W\end{column}
  \begin{column}{0.22\textwidth}\centering\small\textcolor{errorstroke}{\textbf{TinyML}}\\100 mW\end{column}
\end{columns}

\vspace{0.3cm}
{\small\textcolor{lightgray}{What physical constraints drive your decision?}}

\end{frame}

% =============================================================================
\section{Skills \& Prerequisites}
% =============================================================================

\begin{frame}{What You'll Be Able To Do}
\note{[2 min] Concrete, measurable skills. These are the course-level learning
outcomes. Students should feel motivated --- these are real engineering skills.}

\small
By the end of this course, you will be able to:

\vspace{0.1cm}
\begin{enumerate}\setlength\itemsep{0pt}
  \item \textbf{Diagnose bottlenecks} using the Iron Law --- data-bound or compute-bound?
  \item \textbf{Calculate memory budgets} --- can this model fit on this device?
  \item \textbf{Design serving pipelines} with latency and throughput targets
  \item \textbf{Evaluate hardware trade-offs} across the deployment spectrum
  \item \textbf{Compress models} using quantization, pruning, and distillation
  \item \textbf{Reason about training costs} --- time, energy, dollars
  \item \textbf{Monitor production systems} --- detect drift before users do
  \item \textbf{Make principled design decisions} grounded in physics, not hype
\end{enumerate}

\end{frame}

\begin{frame}{Prerequisites}
\note{[2 min] Be honest about what students need. This is a systems course ---
we assume OS and architecture basics. If students lack linear algebra,
recommend a refresher in week 1.}

\small
\begin{columns}[T]
  \begin{column}{0.48\textwidth}
    \textbf{Required background:}

    \vspace{0.1cm}
    \begin{itemize}\setlength\itemsep{0pt}
      \item \textbf{Operating systems} --- processes, memory, I/O
      \item \textbf{Computer architecture} --- caches, pipelines
      \item \textbf{Data structures} --- complexity analysis
      \item \textbf{Linear algebra} --- matrix multiply, norms
      \item \textbf{Probability} --- distributions, Bayes' rule
    \end{itemize}
  \end{column}
  \begin{column}{0.48\textwidth}
    \textbf{Helpful but not required:}

    \vspace{0.1cm}
    \begin{itemize}\setlength\itemsep{0pt}
      \item PyTorch or TensorFlow experience
      \item Systems programming (C/C++)
      \item Cloud computing basics
      \item Prior ML coursework
    \end{itemize}

    \vspace{0.15cm}
    \begin{mlsyscard}{datastroke}
    {\footnotesize \textbf{Note:} We teach ML concepts as needed. No ML expertise required to start.}
    \end{mlsyscard}
  \end{column}
\end{columns}

\end{frame}

\begin{frame}{Course Logistics}
\note{[3 min] Placeholder for instructor to fill in. Cover grading breakdown,
office hours, collaboration policy, late policy, and resources. Adjust per term.}

\small
\begin{columns}[T]
  \begin{column}{0.48\textwidth}
    \textbf{Grading} {\footnotesize(adjust per term)}

    \vspace{0.1cm}
    {\footnotesize
    \renewcommand{\arraystretch}{1.2}
    \begin{tabular}{@{}lr@{}}
      \toprule
      \textbf{Component} & \textbf{Weight} \\
      \midrule
      Problem Sets (5)        & 30\% \\
      Labs (4)                & 20\% \\
      Midterm Exam            & 20\% \\
      Final Project           & 25\% \\
      Participation           & 5\% \\
      \bottomrule
    \end{tabular}
    }
  \end{column}
  \begin{column}{0.48\textwidth}
    \textbf{Resources}

    \vspace{0.1cm}
    \begin{itemize}\setlength\itemsep{1pt}
      \item {\footnotesize \textbf{Textbook}: \textit{ML Systems}, Vol. I}
      \item {\footnotesize \textbf{Office hours}: TBD}
      \item {\footnotesize \textbf{Discussion forum}: TBD}
      \item {\footnotesize \textbf{Hardware labs}: Arduino, Raspberry Pi}
    \end{itemize}

    \vspace{0.15cm}
    \begin{mlsyscard}{routingstroke}
    {\footnotesize \textbf{Collaboration policy:} Discuss concepts freely. Write your own solutions. Cite all sources.}
    \end{mlsyscard}
  \end{column}
\end{columns}

\end{frame}

% =============================================================================
\section{A Taste}
% =============================================================================

\begin{frame}{Why Does GPT-4 Cost \$100M to Train?}
\note{[3 min] Compelling worked example. Walk through the back-of-envelope
calculation. The point: these are not magic numbers --- they follow from
the Iron Law. This is the kind of reasoning you will learn.
Ask: ``Where does the \$100M go?''}

\footnotesize
\textbf{A back-of-envelope calculation:}

\vspace{0.1cm}
\renewcommand{\arraystretch}{1.1}
\begin{tabular}{@{}llr@{}}
  \toprule
  \textbf{Factor} & \textbf{Estimate} & \textbf{Source} \\
  \midrule
  Parameters          & $\sim$1.8 trillion & Architecture choice \\
  Training FLOPs      & $\sim 2 \times 10^{25}$ & $\approx 6 \times N \times D_{\text{tokens}}$ \\
  Hardware            & $\sim$25,000 H100 GPUs & $R_{\text{peak}}$ per chip \\
  Training time       & $\sim$90 days & $O / (R_{\text{peak}} \cdot \eta \cdot N_{\text{GPUs}})$ \\
  Electricity         & $\sim$50 GWh & Power $\times$ time \\
  \midrule
  \textbf{Total cost} & \textbf{$\sim$\$100M} & Hardware + energy + ops \\
  \bottomrule
\end{tabular}

\vspace{0.1cm}
\begin{mlsyscard}{crimson}
{\footnotesize \textbf{Every number traces back to the Iron Law.} The cost follows from physics. This course teaches you to derive these numbers.}
\end{mlsyscard}

\end{frame}

\begin{frame}{What If You Had to Run It on a Phone?}
\note{[2 min] Dramatic contrast. GPT-4 needs 3.6 TB of memory at FP32.
A phone has 8 GB. That is a 450$\times$ gap. Compression, quantization,
distillation --- these are the tools of Part III (Optimize).
Leave this as a teaser --- ``We will solve this problem.''}

\footnotesize
\begin{columns}[T]
  \begin{column}{0.55\textwidth}
    \textbf{GPT-4 at full precision (FP32):}
    Model size $\sim$3.6 TB. Phone RAM: 8 GB. \textcolor{errorstroke}{\textbf{Gap: 450$\times$.}}

    \vspace{0.1cm}
    \textbf{How do we close the gap?}

    \vspace{0.05cm}
    \begin{itemize}\setlength\itemsep{0pt}
      \item \textcolor{computestroke}{\textbf{Quantization}} --- FP32 $\to$ INT4 (8$\times$ smaller)
      \item \textcolor{datastroke}{\textbf{Pruning}} --- remove redundant weights
      \item \textcolor{routingstroke}{\textbf{Distillation}} --- train a smaller model
    \end{itemize}

    \vspace{0.05cm}
    Part III of this course teaches all three.
  \end{column}
  \begin{column}{0.42\textwidth}
    \renewcommand{\arraystretch}{1.1}
    \begin{tabular}{@{}lrr@{}}
      \toprule
      \textbf{Tier} & \textbf{RAM} & \textbf{Gap} \\
      \midrule
      Cloud (H100)   & 80 GB  & 45$\times$ \\
      Edge (Orin)    & 32 GB  & 112$\times$ \\
      Mobile         & 8 GB   & 450$\times$ \\
      TinyML         & 384 KB & 9.4M$\times$ \\
      \bottomrule
    \end{tabular}

    \vspace{0.1cm}
    \begin{mlsyscard}{errorstroke}
    You cannot ``just deploy'' a model. The deployment spectrum forces trade-offs at every tier.
    \end{mlsyscard}
  \end{column}
\end{columns}

\end{frame}


\begin{frame}{Muddiest Point}
\note{[1 min] One-minute paper. Collect responses (physical or digital).
Use responses to open the NEXT lecture with targeted clarification.
Angelo \& Cross (1993) --- powerful diagnostic tool.}

\centering
\Large\bfseries One-minute paper\\[0.5cm]
\normalsize Write down the \textbf{muddiest point} from today's lecture.\\[0.2cm]
{\small What concept was most confusing or unclear?}\\[0.5cm]
{\footnotesize\textcolor{midgray}{(Hand in on your way out --- or submit digitally)}}
\end{frame}

% --- ACTIVE LEARNING 3: Retrieval Practice ---
\begin{frame}{What Were the Key Ideas?}
\note{[2 min] Retrieval practice. Students write for 90 seconds, no notes.
Walk around the room. Do NOT show the next slide yet.}

\centering
\vspace{1.5cm}
{\Large\bfseries Close your notes.}

\vspace{0.8cm}
{\large Write down the \textbf{4 most important ideas} from today.}

\vspace{0.8cm}
{\normalsize\textcolor{midgray}{90 seconds --- no peeking.}}

\end{frame}

\begin{frame}{Key Takeaways}
\note{[2 min] Reveal. Walk through each bullet. Emphasize that every idea
introduced today will recur throughout the semester.}

\footnotesize
\begin{itemize}\setlength\itemsep{0pt}
  \item \textbf{Constraints drive architecture.} AI is infrastructure with physical laws --- not magic.
  \item \textbf{The 5\% problem:} Model code is 5\% of production. The other 95\% is systems engineering.
  \item \textbf{\DAM{} Taxonomy:} Data, Algorithm, Machine are interdependent. ``Which is the bottleneck?''
  \item \textbf{Iron Law:} $T = D/BW + O/(R_p \cdot \eta) + L$. Slowest term dominates.
  \item \textbf{Degradation Eq.:} $A(t) = A_0 - \alpha \cdot \Delta(P)$. Models fail silently as distributions drift.
  \item \textbf{Five Lighthouses:} ResNet-50, GPT-2, MobileNetV2, DLRM, KWS --- each isolates a bottleneck.
  \item \textbf{Deployment Spectrum:} Cloud to TinyML spans $10^6\!\times$ compute. Different engineering per tier.
  \item \textbf{Course structure:} Foundations $\to$ Build $\to$ Optimize $\to$ Deploy across 16 chapters.
\end{itemize}

\end{frame}

\begin{frame}{Next Lecture: Chapter 1 --- Introduction}
\note{[1 min] Forward hook. Chapter 1 goes deep on the D-A-M taxonomy,
the Iron Law with worked examples, and the five lighthouse models.
Assign pre-reading if applicable.}

\small
\centering
\vspace{0.5cm}
{\large\textbf{Chapter 1: Introduction}}

\vspace{0.3cm}
{\normalsize The Physics of AI Engineering}

\vspace{0.5cm}
\begin{columns}[T]
  \begin{column}{0.30\textwidth}
    \centering
    \textcolor{datastroke}{\textbf{\DAM{} Taxonomy}}\\[0.1cm]
    {\footnotesize Deep dive into all three axes}
  \end{column}
  \begin{column}{0.30\textwidth}
    \centering
    \textcolor{computestroke}{\textbf{Iron Law}}\\[0.1cm]
    {\footnotesize Worked examples with real hardware}
  \end{column}
  \begin{column}{0.30\textwidth}
    \centering
    \textcolor{errorstroke}{\textbf{Degradation}}\\[0.1cm]
    {\footnotesize Why models fail silently}
  \end{column}
\end{columns}

\vspace{0.5cm}
{\normalsize\textbf{Read:} Chapter 1 of the textbook before next class.}

\end{frame}


% =============================================================================
% BACKUP SLIDES
% =============================================================================
\appendix

\begin{frame}{Backup: Iron Law Dimensional Analysis}
\note{Use for additional depth or if students need alternative explanation.}
\small
Use if students ask about dimensional consistency.

\small
Each Iron Law term resolves to \textbf{seconds}:

\vspace{0.2cm}
\renewcommand{\arraystretch}{1.2}
{\footnotesize
\begin{tabular}{@{}lll@{}}
  \toprule
  \textbf{Term} & \textbf{Units} & \textbf{Resolves to} \\
  \midrule
  $D_{\text{vol}} / BW$ & Bytes / (Bytes/s) & seconds \\
  $O / (R_p \cdot \eta)$ & FLOPs / (FLOPs/s) & seconds \\
  $L_{\text{lat}}$ & seconds & seconds \\
  \bottomrule
\end{tabular}
}

\vspace{0.2cm}
You can always check your calculation: if the units do not resolve to seconds, the equation is wrong.
\end{frame}

\begin{frame}{Backup: Deployment Spectrum Extended}
\note{Use for additional depth or if students need alternative explanation.}
\small
Use if students want more detail on TinyML vs.\ Mobile.

\small
\renewcommand{\arraystretch}{1.15}
{\footnotesize
\begin{tabular}{@{}lllll@{}}
  \toprule
  \textbf{Tier} & \textbf{Device} & \textbf{RAM} & \textbf{Compute} & \textbf{Power} \\
  \midrule
  Cloud & H100 & 80 GB HBM3 & 989 TF & 700 W \\
  Edge & Orin NX & 32 GB LPDDR5 & 100 TOPS & 15--40 W \\
  Mobile & Flagship SoC & 8 GB LPDDR5 & 17 TOPS & 3--5 W \\
  TinyML & Arduino Nicla & 384 KB SRAM & 0.5 GF & 100 mW \\
  \bottomrule
\end{tabular}
}

\vspace{0.15cm}
The gap from Cloud to TinyML spans $10^5\!\times$ in memory and $10^6\!\times$ in compute. Each tier requires fundamentally different engineering.
\end{frame}

\end{document}
